\chapter{图论}
\label{chap:graph-theory}
第~\ref{chap:combinatorics}章用\(z_1,\ldots,z_4\)四个二值变量说明有限变量、约束
与目标。本章不增加变量，而是直接从式~\eqref{eq:comb-running-cnf}的约束作用域构造
\term{原始关系图}（primal graph）：每个二值变量对应一个顶点，同一子句中的任意两个
变量相连。前两个子句给出边\(\{1,2\}\)，第三个子句给出边\(\{1,3\}\)，三变量
子句\(\neg z_4\lor z_2\lor z_3\)则使\(2,3,4\)三个顶点两两相连；这样的顶点集合
后文称为团。图只记录哪些变量需要在同一个局部约束中共同处理；三变量子句仍是一个三元
约束，不能把这三条边误解为三个与原子句等价的二元约束。

贯穿例固定第~\ref{chap:combinatorics}章的四个二值变量：
\[
  z_i\in\{0,1\},\qquad i\in V=\{1,2,3,4\}.
\]
讨论不依赖二值性的图算法时，用有限集合\(\mathcal{Z}_i\)表示变量\(z_i\)的一般
取值域；贯穿例对应\(\mathcal{Z}_i=\{0,1\}\)。
图~\ref{fig:graph-theory-running-graph}中的边表示哪些变量共同出现在某个约束中。
边集合为
\begin{equation}
  E=\bigl\{\{1,2\},\{1,3\},\{2,3\},\{2,4\},\{3,4\}\bigr\}.
  \label{eq:graph-theory-running-edges}
\end{equation}
这张图很小，却包含本章需要分开的几种作用。两个三角形\(\{1,2,3\}\)与
\(\{2,3,4\}\)共享接口\(\{2,3\}\)，所以固定接口后，只含\(z_1\)与只含\(z_4\)
的局部计算可以分开。若过早消去顶点\(2\)，原本不相邻的顶点\(1,4\)会在中间结果中
直接联系。同一连接结构可以承载硬约束、概率、得分、距离或容量，但这些数值语义并不
由图本身决定。

\begin{figure}[htbp]
  \centering
  \begin{tikzpicture}[
    vertex/.style={circle,draw=BookTeal,fill=BookTeal!8,
      minimum size=8mm,inner sep=0pt,font=\figurefont},
    edge/.style={draw=BookMuted,line width=0.9pt}
  ]
    \node[vertex] (x2) at (0,1.15) {\(z_2\)};
    \node[vertex] (x3) at (0,-1.15) {\(z_3\)};
    \node[vertex] (x1) at (-2.1,0) {\(z_1\)};
    \node[vertex] (x4) at (2.1,0) {\(z_4\)};
    \draw[edge] (x1)--(x2);
    \draw[edge] (x1)--(x3);
    \draw[edge] (x2)--(x3);
    \draw[edge] (x2)--(x4);
    \draw[edge] (x3)--(x4);
  \end{tikzpicture}
  \caption{由第~\ref{chap:combinatorics}章约束作用域导出的变量关系图。每条边表示
  两个变量共同出现在至少一个约束中；共享边\(\{2,3\}\)是左右两个三角形的接口。
  边权、方向与关系的数值含义需要另行给出。}
  \label{fig:graph-theory-running-graph}
\end{figure}

本章先建立图的对象与表示，再研究遍历、分离和树分解怎样揭示局部结构。分解只说明
计算应如何组织，真正计算什么还由代数运算决定，因此第
~\ref{sec:graph-theory-semiring-dynamic-programming}节用半环统一边缘化、最优赋值
与最短路径。第~\ref{sec:graph-theory-exact-approximate-algorithms}节比较几类精确图算法
及其条件，并完整推导二值次模能量的最小割表示；
一般离散目标不具备这些结构时，仍可用枚举或分支定界精确搜索，但最坏情形通常代价很高；
也可借助松弛、取整或局部搜索获得近似解与上下界证书。最后一节把顶点上的数值看作图信号，
用图拉普拉斯连接平滑、分割与扩散，并检验顶点重编号下哪些表示保持不变。

\section{图的对象与表示}
\label{sec:graph-theory-objects-representations}

\subsection{顶点与边只规定关系的骨架}

\term{无向图}（undirected graph）是二元组
\[
  \mathcal{G}=(V,E),
\]
其中\(V\)是顶点集合，\(E\)是由不同顶点组成的无序对集合。若
\(\{i,j\}\in E\)，就称\(i,j\)相邻，并把二者分别称为对方的邻居。顶点\(i\)的
\term{邻域}（neighborhood）与\term{度}（degree）分别为
\[
  N(i)=\{j\in V:\{i,j\}\in E\},\qquad d_i=|N(i)|.
\]
式~\eqref{eq:graph-theory-running-edges}定义的图中，
\(N(2)=\{1,3,4\}\)，而\(N(1)=\{2,3\}\)。这里采用不含自环、平行边的有限简单图；
若任务允许自环或同一对顶点间有多条边，必须把边集合改为相应的多重结构。

无向边表达对称关系。若关系具有先后或作用方向，就使用
\term{有向图}（directed graph）\(\mathcal{G}=(V,A)\)，其中
\(A\subseteq V\times V\)由有序对组成。\((i,j)\in A\)表示从\(i\)指向\(j\)，
不自动推出\((j,i)\in A\)。指向顶点\(j\)的弧称为\(j\)的入边，入边条数称为
\term{入度}（in-degree）；从\(j\)指出的弧称为出边，其条数称为
\term{出度}（out-degree）。因果生成次序、状态转移和计算依赖通常需要有向边；
距离邻接或对称相似性通常使用无向边。

边还可以带有\term{权重}（weight）。无向加权图要求
\(w_{i,j}=w_{j,i}\)，有向加权图则允许二者不同。权重可以表示长度、容量、相似度
或局部代价；这些量对“权重大”的解释不同，所以每个算法都须重新声明权重语义。本章
不统一排除零权边：结构边集\(E\)记录哪些关系存在，而对无向非负权图，
\[
  E_+=\bigl\{\{i,j\}\in E:w_{i,j}>0\bigr\}
\]
称为\term{正权支撑}（positive-weight support）。当权重作为线性算子的系数时，
零权结构边不产生贡献，相关的谱与扩散结论因而取决于支撑图\(\mathcal{G}_+=(V,E_+)\)。

成对边有时仍不足以表达局部关系。若一个约束同时作用于\(\{1,3,4\}\)，把它拆成三条
边可能改变允许的赋值。此时使用\term{超图}（hypergraph），其超边可以包含两个以上
顶点。也可使用\term{二部图}（bipartite graph）把顶点分成两组，只允许跨组边。
变量与局部因子形成的因子图就是二部图，它保留了高阶关系同时涉及哪些变量。

\subsection{路径与环描述局部边怎样组成整体连接}

图中的\term{游走}（walk）是顶点序列
\(v_0,v_1,\ldots,v_k\)，相邻两项之间有边；顶点与边可以重复。若
\(v_0,\ldots,v_k\)两两不同，就得到\term{路径}（path）。长度是经过的边数；
在带长度的图上，路径长度改为边长之和。
无向图中的\term{环}（cycle）是长度\(k\geq3\)的闭游走
\(v_0,v_1,\ldots,v_k\)，其中\(v_0=v_k\)，而\(v_0,\ldots,v_{k-1}\)两两不同。
因此环只在首尾重复同一顶点，不与路径要求全部顶点互异的定义混用。

若任意两个顶点间都有路径，无向图称为\term{连通}（connected）。一般无向图可以唯一
分成若干极大的连通子图，称为连通分量。图
~\ref{fig:graph-theory-running-graph}连通，因为每个顶点都能经顶点\(3\)到达其他顶点；
\(1,2,3,1\)与\(2,3,4,2\)分别构成一个环。删除边\(\{1,2\}\)后图仍连通，
但第一个环消失。

连通且无环的无向图称为\term{树}（tree）。树中任意两点之间恰有一条简单路径：
两条不同路径从分开到重新会合的部分会构成环。含\(n\)个顶点的树恰有\(n-1\)条边，
若干棵互不相连的树组成森林。下面把路径唯一性和边数结论合在一起证明，因为二者分别
解释树为何适合递推，以及递推中为何不会遗漏额外连接。

\begin{theorem}[树的等价刻画]
  \label{thm:graph-theory-tree-characterizations}
  对含\(n\geq1\)个顶点的有限无向简单图，以下三条等价：
  \begin{enumerate}
    \item 图连通且无环；
    \item 任意两点之间恰有一条简单路径；
    \item 图连通，并且恰有\(n-1\)条边。
  \end{enumerate}
\end{theorem}

\begin{proof}
若图连通且无环，任意两点间至少有一条简单路径；若有两条不同的简单路径，从二者首次
分开处走到再次会合处便得到一个环，产生矛盾，故路径唯一。反过来，路径唯一已经包含
连通性；若图有环，环上相邻顶点之间既有该边，又有沿环另一侧的路径，与唯一性矛盾。

再证明边数。单顶点树有\(0\)条边。含至少两个顶点的有限树必有叶顶点：取一条最长
简单路径，端点若还有路径外邻居，就能延长该路径；端点若连接路径中的非相邻顶点，
又会形成环。删除叶顶点及其唯一关联边后仍是树，归纳得到边数
\((n-1)-1+1=n-1\)。树中任意边又位于其两端的唯一简单路径上，删除该边必使两端
断开。

最后，设图连通且有\(n-1\)条边。从任一顶点开始，每次加入一条连接“已加入顶点”与
“尚未加入顶点”的边；连通性保证只要还有顶点未加入，这样的边就存在。加入一个新顶点
不会形成环，重复\(n-1\)次便得到一棵覆盖全部顶点的树。该树已经使用原图的全部
\(n-1\)条边，所以原图与这棵树相同，因而无环。
\end{proof}

有向图中的路径必须顺着弧方向前进。\term{有向环}（directed cycle）是除首尾外顶点
互异、每一步都沿弧方向的闭游走。在不含自环的有向图中，有向环长度至少为\(2\)：
若\(i\to j\)与\(j\to i\)同时存在，二者已经构成有向环；这与无向简单图的环至少
包含三个顶点不同。不含有向环的图称为
\term{有向无环图}（directed acyclic graph, DAG）。DAG中的边可以规定偏序依赖，
但图中没有边的两个顶点不一定可以互换执行，它们还可能通过更长路径间接依赖。
后文的拓扑排序会把这种偏序扩展为一个线性次序。

有向图还需区分两种连通性。若任意有序顶点对\(i,j\)之间都有从\(i\)到\(j\)的
有向路径，图称为\term{强连通}（strongly connected）；若忽略弧方向后所得无向图
连通，则称为\term{弱连通}（weakly connected）。强连通蕴含弱连通，反向不成立。
例如单向链\(1\to2\to3\)弱连通，却不存在从\(3\)返回\(1\)的有向路径。

\subsection{邻接矩阵把关系写成线性映射}

给顶点固定次序\(1,\ldots,n\)，无向图的\term{邻接矩阵}（adjacency matrix）
\(\bm{A}\in\mathbb{R}^{n\times n}\)定义为
\[
  a_{i,j}=\begin{cases}
    w_{i,j},&\{i,j\}\in E,\\
    0,&\text{其他情形}.
  \end{cases}
\]
无权图取\(w_{i,j}=1\)。无向图使\(\bm{A}\)对称；有向图则令
\(a_{i,j}\)记录弧\(i\to j\)的权重，一般不对称。简单无向图没有自环，所以主对角线
为零。若允许结构边的权重为\(0\)，数值矩阵中的零条目不能区分零权边与无边，此时必须
把\(\bm{A}\)与边集\(E\)或一个邻接掩码一同保存；只有在所有结构边权非零时，
\(\bm{A}\)才单独确定边集。图~\ref{fig:graph-theory-running-graph}的邻接矩阵是
\begin{equation}
  \bm{A}=\begin{bmatrix}
    0&1&1&0\\
    1&0&1&1\\
    1&1&0&1\\
    0&1&1&0
  \end{bmatrix}.
  \label{eq:graph-theory-running-adjacency}
\end{equation}

邻接矩阵不只是存储表。对顶点数值
\(\bm{f}=(f_1,\ldots,f_n)^{\top}\)，
\[
  (\bm{A}\bm{f})_i=\sum_{j\in N(i)}w_{i,j}f_j
\]
汇总顶点\(i\)的邻居值，因此是第~\ref{chap:linear-algebra}章意义下的线性映射。
矩阵幂还记录游走：对无权图，
\((\bm{A}^k)_{i,j}\)等于从\(i\)到\(j\)长度为\(k\)的游走数。证明只需对\(k\)
归纳；矩阵乘法在中间顶点\(\ell\)上求和，恰好把长度\(k\)的游走与最后一条边拼接。
游走允许重复顶点，所以该数一般不等于简单路径数。对带权图，同一归纳说明
\[
  (\bm{A}^k)_{i,j}
  =\sum_{i=v_0\to v_1\to\cdots\to v_k=j}
  \prod_{r=1}^{k}w_{v_{r-1},v_r},
\]
即条目是所有长度\(k\)游走的权重乘积之和，而不再是游走条数。这个解释依赖通常的
加法与乘法；换用布尔或最小加法运算时，矩阵幂会回答可达性或固定边数最短成本等不同
问题，不能继续称为计数。

矩阵表示依赖顶点次序。若置换矩阵\(\bm{P}\)改变编号，同一无向图在新编号下表示为
\begin{equation}
  \bm{A}'=\bm{P}\bm{A}\bm{P}^{\top}.
  \label{eq:graph-theory-adjacency-relabel}
\end{equation}
矩阵元素的位置改变，连通、环和路径等图性质不应改变。反过来，并非任意相似变换都只是
顶点重编号；图重编号只允许置换矩阵，而第~\ref{chap:matrix-analysis}章中的一般换基
会混合顶点坐标，通常不再得到同一张离散图。

还可以给每条无向边任取一个方向，构造\term{关联矩阵}（incidence matrix）
\(\bm{B}\in\mathbb{R}^{|V|\times|E|}\)：边\(e:i\to j\)对应列在\(i\)处取\(1\)、
在\(j\)处取\(-1\)，其余为零。反转临时方向只会把该列乘\(-1\)，不改变后文的
\(\bm{B}\bm{W}\bm{B}^{\top}\)。邻接矩阵为每对顶点预留位置，适合稠密线性运算；
邻接表只存每个顶点实际相邻的对象，空间为\(O(|V|+|E|)\)，更适合稀疏图遍历；
关联矩阵则显式保留“哪条边连接哪些顶点”，便于表达边差分。在结构边权均非零，或为
邻接矩阵另附边集或掩码时，三种表示保留同一个图对象，但直接支持的计算和存储成本不同。

\subsection{图对象与赋予图的语义必须分开}

在贯穿关系图上，可以为每个顶点放置一元函数
\(\phi_i(z_i)\)，并在只考虑一元与成对局部项时为每条边放置成对函数
\(\phi_{i,j}(z_i,z_j)\)，形成
\begin{equation}
  F(\bm{z})=
  \bigotimes_{i\in V}\phi_i(z_i)\otimes
  \bigotimes_{\{i,j\}\in E}\phi_{i,j}(z_i,z_j).
  \label{eq:graph-theory-factorized-objective}
\end{equation}
符号\(\otimes\)暂时只表示组合局部信息，第
~\ref{sec:graph-theory-semiring-dynamic-programming}节才规定它可能是乘法或加法。
若要继续表示式~\eqref{eq:comb-running-cnf}的原始硬约束，第四个子句必须保留为
三元因子\(\phi_{2,3,4}(z_2,z_3,z_4)\)；式
~\eqref{eq:graph-theory-factorized-objective}是同一关系骨架上的成对专门情形，
不是把三元子句拆成三条边后的等价公式。
图说明每个局部项的作用域，却不说明这些函数是否为概率、约束指示、相似度或成本。
同一张图既可以计算总权重，也可以寻找最低成本配置；反过来，相同数值表若采用不同
作用域，也会产生不同的分解与计算代价。

图不会自动编码条件独立、因果或时间关系；这些语义还需生成规则、分布因子或结构假设。

\section{遍历、分离与结构分解}
\label{sec:graph-theory-traversal-separation-decomposition}

\subsection{遍历把局部邻接组织成可执行次序}

\term{广度优先搜索}（breadth-first search, BFS）从起点\(s\)出发，先访问距离
\(s\)为一条边的顶点，再访问距离为两条边的顶点。维护一个先进先出队列，并在顶点首次
入队时标记它，可以保证每个顶点至多入队一次、每条边至多检查常数次。因此使用邻接表时，
时间复杂度为\(O(|V|+|E|)\)\scite{math-cormen2022algorithms}每个非起点顶点首次
被发现时记录一条父边，这些父边不成环，并在所有可达顶点上构成一棵搜索树；队列、已访问
集合和父边分别回答“接下来处理谁”“谁已处理过”和“如何恢复发现路径”。

对无权图，BFS首次发现顶点\(v\)时记录的层数就是从\(s\)到\(v\)的最短路径长度。
证明按队列层次归纳：处理第\(k\)层顶点时，所有更短路径的终点已经处理；新发现邻居有
一条长度\(k+1\)的路径。若它存在更短路径，其倒数第二个顶点应位于更早层，并会更早
发现它，产生矛盾。带一般边长时，边数最少不等于总长度最小，BFS便不再适用。

\term{深度优先搜索}（depth-first search, DFS）沿一条尚未探索的边持续深入，无法
继续时回退。它同样能在线性时间内找连通分量，还能依据进入与退出时间识别有向图中的
回边。若DFS遇到一条指向当前递归栈内顶点的弧，就得到一个有向环；若没有这种回边，
按退出时间逆序排列顶点便得到拓扑序。

拓扑序与有向无环性的等价关系是图论中的基本结构结论\scite{math-diestel2017graph}
\begin{theorem}[DAG的拓扑排序]
  \label{thm:graph-theory-topological-order}
  有限有向图存在拓扑序，当且仅当它不含有向环。拓扑序是顶点的一个排列，使每条弧
  \(i\to j\)中\(i\)都排在\(j\)之前。
\end{theorem}

\begin{proof}
若存在有向环
\(v_1\to v_2\to\cdots\to v_k\to v_1\)，拓扑序将依次要求
\(v_1\)在\(v_2\)前、\(v_2\)在\(v_3\)前，最终又要求\(v_k\)在\(v_1\)前，
不可能同时满足。

反过来，设图无有向环。图中必有入度为零的顶点；否则从任意顶点\(v\)出发，总能选择
当前顶点的一条入边\(u\to v\)，并反向追溯到其起点\(u\)。顶点数有限，追溯中必有
顶点重复，而重复段按原弧方向组成有向环，产生矛盾。删除一个入度为零的顶点及其出边，
剩余图仍无环。重复这一过程直到图为空，删除次序就是拓扑序。每次若用队列维护新出现的
零入度顶点，全部过程只检查每个顶点和每条弧常数次。
\end{proof}

拓扑序一般不唯一。两个互相不可达的顶点可以有多种相对次序；这反映偏序未规定它们的
先后，而不是算法给出了矛盾答案。

\subsection{分隔集把全局连接切成局部部分}

对无向图，若删除顶点集合\(S\)及其关联边后，顶点集合\(A\)与\(B\)落入不同连通分量，
就称\(S\)将\(A\)与\(B\)\term{分隔}（separate）。等价地，从\(A\)中任一点到
\(B\)中任一点的每条路径都经过\(S\)。在贯穿图中，集合\(\{2,3\}\)分隔
\(\{1\}\)与\(\{4\}\)。一旦固定\(z_2,z_3\)的接口信息，只含\(z_1\)与只含
\(z_4\)的局部项就可以分别处理，再通过接口组合。

删除一个顶点及其关联边后使连通分量数增加的顶点称为割点；删除一条边后使连通分量数
增加的边称为桥。贯穿图的每条边都位于某个三角形中，删除该边后仍有连接其两端的替代
路径，所以没有桥。删除顶点\(1\)或\(4\)后，余图是一个三角形；删除顶点\(2\)或
\(3\)后，余图分别是路径\(1-3-4\)或\(1-2-4\)，仍然连通，所以也没有割点。
但删除两个顶点组成的分隔集\(\{2,3\}\)会使\(1\)与\(4\)断开。桥描述连接冗余，
不能单凭边权大小判断。

两两相邻的顶点集合称为\term{团}（clique）。若它不被更大的团包含，就称为极大团。
贯穿图的极大团为
\[
  \{1,2,3\},\qquad \{2,3,4\}.
\]
“极大”表示不能再加入顶点，“最大”则表示顶点数在全部团中最多；两者不应混用。
团把所有成对关系集中到同一个局部集合，但大团也意味着局部表需要同时索引更多变量。

\subsection{消元会制造填充边}

要从式~\eqref{eq:graph-theory-factorized-objective}中消去变量\(z_i\)，必须先组合
所有含\(z_i\)的局部项，再对\(z_i\)汇总。结果通常成为\(N(i)\)上一个新的联合函数，
所以原本互不相邻的邻居会在中间结果中直接联系。图上把这些邻居补成团，并删除\(i\)，
称为一次\term{消元}（elimination）；新加入的边称为\term{填充边}（fill edge）。

对贯穿图采用次序
\[
  \pi_{\mathrm{good}}=(1,4,2,3).
\]
消去\(1\)时，其尚未消去的邻居\(\{2,3\}\)已经相邻；消去\(4\)时同理。
整个过程中，任何顶点被消去时至多有两个尚未消去的邻居，因此不会产生填充边。

若采用
\[
  \pi_{\mathrm{bad}}=(2,1,3,4),
\]
第一步就要把邻居\(\{1,3,4\}\)补成团，加入原图中不存在的边\(\{1,4\}\)。
图的原始边数没有改变问题的候选空间，但消元次序改变了中间函数必须同时依赖的变量数。

给定消元次序，某一步中尚未消去邻居的最大数量称为该次序的
\term{诱导宽度}（induced width）。上述两个次序的诱导宽度分别为\(2\)和\(3\)。
若每个变量至多有\(q\)个取值，作用于\(k+1\)个变量的稠密局部表最多含
\(q^{k+1}\)项，因此宽度进入指数，而不仅是一个图形美观指标。

\subsection{弦图用局部消元刻画无长无弦环}

长度至少为\(4\)的环若有一条连接非相邻环顶点的边，这条边称为弦。每个这类环都有弦
的图称为\term{弦图}（chordal graph）。树因没有环而是弦图；四顶点环加入任一对角线
后才成为弦图。

若图中任意两个不同顶点都相邻，就称该图为\term{完全图}（complete graph）。给定
顶点集\(U\subseteq V\)，保留\(U\)中的顶点以及原图中两个端点都在\(U\)内的全部
边，所得图称为\(U\)上的\term{诱导子图}（induced subgraph），记作
\(\mathcal{G}[U]\)。

若一个顶点的邻居组成团，就称它为\term{单纯顶点}（simplicial vertex）。
按某个次序依次删除顶点时，若每个顶点在剩余图中都是单纯顶点，该次序称为
\term{完美消元序}（perfect elimination ordering）。

\begin{lemma}[弦图的单纯顶点]
  \label{lem:graph-theory-chordal-simplicial-vertex}
  每个非空有限弦图都有单纯顶点；若图不是完全图，则有两个互不相邻的单纯顶点。
\end{lemma}

\begin{proof}
对顶点数归纳。完全图中每个顶点都单纯。若图不连通，对两个连通分量分别应用归纳假设，
各取一个单纯顶点；不同分量的两个顶点不相邻，并且它们在分量内与全图中的邻域相同。

下面设图连通但不完全。取两个不相邻顶点\(a,b\)，并取一个按集合包含关系极小的
\(a\)--\(b\)分隔集\(S\)，即对每个\(s\in S\)，集合\(S\setminus\{s\}\)
都不再分隔\(a,b\)。先说明\(S\)是团。记删除\(S\)后包含\(a,b\)的分量分别为
\(C_a,C_b\)。这种极小性保证每个\(s\in S\)在\(C_a,C_b\)中都各有邻居，否则可以
从\(S\)删去\(s\)而仍分隔\(a,b\)。若\(x,y\in S\)不相邻，分别在
\(C_a\cup\{x,y\}\)与\(C_b\cup\{x,y\}\)中取从\(x\)到\(y\)的最短路径。
两条路径合成长度至少为\(4\)的环；路径最短性排除同侧的弦，两个分量之间没有边，
而\(x,y\)也不相邻，故该环无弦，产生矛盾。因此\(S\)是团。

对任一\(C\in\{C_a,C_b\}\)，诱导子图\(\mathcal{G}[C\cup S]\)是顶点更少的
弦图。若它是完全图，任取\(C\)中顶点即为该诱导子图的单纯顶点；若它不完全，归纳假设
给出两个互不相邻的单纯顶点，而\(S\)是团，所以至少一个位于\(C\)。取这个位于\(C\)
的顶点\(v_C\)。顶点\(v_C\)在原图中没有\(C\cup S\)之外的邻居，故仍是原图的
单纯顶点。分别从\(C_a,C_b\)取得的两个顶点互不相邻，结论成立。
\end{proof}

\begin{theorem}[弦图与完美消元序]
  \label{thm:graph-theory-chordal-perfect-elimination}
  有限无向图是弦图，当且仅当它具有完美消元序。
\end{theorem}

\begin{proof}
若图为弦图，引理~\ref{lem:graph-theory-chordal-simplicial-vertex}给出一个单纯顶点。
删除该顶点得到的诱导子图仍为弦图，故可重复应用引理；依次删除的顶点组成完美消元序。

反过来，设图有完美消元序但含一个长度至少为\(4\)的无弦环。取环上最早被消去的顶点
\(v\)。此时它的两个环上邻居都尚未消去；无弦性说明这两个邻居不相邻，所以\(v\)在
剩余图中不是单纯顶点，与完美消元序的定义矛盾。
\end{proof}

这个判据还给出消元接口：任一消元次序加入全部填充边后都会得到弦图，而该次序在所得
弦图中成为完美消元序。

贯穿图本来就是弦图。次序\(\pi_{\mathrm{good}}\)是完美消元序，因为顶点\(1,4\)
各自的剩余邻域都是团\(\{2,3\}\)。次序\(\pi_{\mathrm{bad}}\)不是原图的完美消元序；
补上边\(\{1,4\}\)后，它才成为完成图的完美消元序。弦图不是说消元计算代价一定小，
大团仍可使宽度很大。

\subsection{树分解把一般图组织成相互一致的局部袋}

\term{树分解}（tree decomposition）由一棵树\(\mathcal{T}\)及每个树顶点
\(t\)对应的顶点袋\(B_t\subseteq V\)组成，并满足：
\begin{enumerate}
  \item 每个原图顶点至少出现在一个袋中；
  \item 每条原图边的两个端点共同出现在某个袋中；
  \item 对任意顶点\(i\)，所有满足\(i\in B_t\)的树顶点在\(\mathcal{T}\)中连通。
\end{enumerate}
第三条称为运行交集条件。它保证一个变量若出现在多个局部袋中，这些副本可以沿树上的
连续路径保持一致；仅覆盖全部边而不满足这条条件，不能支持可靠的局部消息传递。

\begin{lemma}[树上连通子树的 Helly 性质]
  \label{lem:graph-theory-subtree-helly}
  设\(\mathcal{T}_1,\ldots,\mathcal{T}_m\)是有限树\(\mathcal{T}\)中的非空连通
  子树。若任意两棵子树都有公共顶点，则全部子树有公共顶点：
  \[
    \mathcal{T}_i\cap\mathcal{T}_j\neq\varnothing
    \quad(1\leq i,j\leq m)
    \quad\Longrightarrow\quad
    \bigcap_{i=1}^{m}\mathcal{T}_i\neq\varnothing.
  \]
\end{lemma}

\begin{proof}
在\(\mathcal{T}\)中任选根。对每棵\(\mathcal{T}_i\)，令\(a_i\)为其中离根最近的
顶点。若最近顶点不唯一，或\(a_i\)不在根到某个
\(x\in\mathcal{T}_i\)的路径上，则连接这两个子树顶点的唯一路径会经过一个离根更近
的顶点；连通子树包含这条路径，产生矛盾。因此\(a_i\)唯一，并且位于根到
\(\mathcal{T}_i\)中每个顶点的路径上。取深度最大的\(a_k\)。

固定任意\(i\)。由两两相交，可取
\(x\in\mathcal{T}_i\cap\mathcal{T}_k\)。顶点\(a_i,a_k\)都在从根到\(x\)的
路径上，而\(a_k\)的深度不小于\(a_i\)，所以\(a_k\)位于\(a_i\)到\(x\)的路径上。
\(\mathcal{T}_i\)连通且含\(a_i,x\)，必包含二者之间的整条路径，因而
\(a_k\in\mathcal{T}_i\)。这对每个\(i\)都成立，故\(a_k\)属于全部子树。
\end{proof}

这个引理补上了团与袋之间的关键一步。对原图顶点\(v\)，记
\(\mathcal{T}_v=\{t:v\in B_t\}\)。运行交集条件说明\(\mathcal{T}_v\)是连通
子树；若\(u,v\)相邻，边覆盖条件又说明
\(\mathcal{T}_u\cap\mathcal{T}_v\neq\varnothing\)。因此，对任意团\(C\)，
\(\{\mathcal{T}_v:v\in C\}\)两两相交。由引理
~\ref{lem:graph-theory-subtree-helly}，存在同一个树顶点\(t\)属于所有这些子树，
也就是\(C\subseteq B_t\)。所以大小为\(r\)的团给出树宽下界\(r-1\)。

贯穿图有树分解
\[
  B_a=\{1,2,3\},\qquad B_b=\{2,3,4\},
\]
其中让\(B_a-B_b\)连接成一条树。顶点\(2,3\)各出现在两个袋中且这些袋连通，
顶点\(1,4\)各出现一次。每条边也被某个袋覆盖。

树分解的宽度定义为
\[
  \max_{t}|B_t|-1.
\]
减一使树的树宽为\(1\)，单个顶点图的树宽为\(0\)。图的
\term{树宽}（treewidth）是所有树分解宽度的最小值。贯穿图含三角形；刚才的团下界
说明任何树分解都必须在某个袋中同时容纳三角形的三个顶点，所以树宽至少为\(2\)。
上述分解的两个袋大小均为\(3\)，宽度达到\(2\)，故贯穿图的树宽恰为\(2\)。

\begin{theorem}[消元序产生树分解]
  \label{thm:graph-theory-elimination-tree-decomposition}
  给定有限无向图及一个完整消元次序。若顶点\(v\)被消去时尚未消去的邻居集合为
  \(N^+_{\pi}(v)\)，并在消元时把它们补成团，则存在袋
  \[
    B_v=\{v\}\cup N^+_{\pi}(v)
  \]
  组成的树分解，其宽度不超过该次序的诱导宽度。因而图的树宽不大于任一消元次序的
  诱导宽度。
\end{theorem}

\begin{proof}
在完成消元后得到的弦图中，每个\(N^+_{\pi}(v)\)都是团。对每个非空
\(N^+_{\pi}(v)\)，令\(p(v)\)是该集合中最早被消去的顶点，并连接袋\(B_v\)与
\(B_{p(v)}\)；若集合为空，则把对应袋作为一个根。若原图不连通，可以把各棵所得树
任意连接，空交集不会破坏条件。

每条连接都从较早消去顶点指向较晚消去顶点，所以不可能形成有向环；除根外每个袋恰有
一个父袋，因而得到树。每个顶点\(v\)属于自己的袋。若原图边\(\{u,v\}\)中\(u\)
先被消去，则\(v\in N^+_{\pi}(u)\)，所以该边被\(B_u\)覆盖。

还需验证运行交集。固定顶点\(z\)。若\(z\in B_v\)且\(v\neq z\)，则
\(z\in N^+_{\pi}(v)\)。由于该集合已补成团，\(z\)与\(p(v)\)相邻；若
\(p(v)\neq z\)，则\(z\in B_{p(v)}\)。因此从任一含\(z\)的袋沿父边向上走，在到达
\(B_z\)之前经过的每个袋都仍含\(z\)。全部含\(z\)的袋于是连成子树。

最后，\(|B_v|-1=|N^+_{\pi}(v)|\)，最大袋宽不超过该消元次序的诱导宽度。
\end{proof}

反方向也成立：宽度为\(k\)的树分解可通过依次删除叶袋中特有顶点，构造诱导宽度不超过
\(k\)的消元次序。证明维持如下不变量：每轮开始时，剩余袋构成当前填充图的树分解，而
不只覆盖原图。初始时这正是给定树分解。先删除被相邻袋包含的冗余叶袋；对剩余叶袋
\(B_\ell\)及其父袋\(B_p\)，集合\(X=B_\ell\setminus B_p\)非空。运行交集条件说明
\(X\)中的顶点只出现在\(B_\ell\)，再由当前填充图的边覆盖条件可知，这些顶点的尚存
邻居都位于\(B_\ell\)。

依次消去\(X\)中的顶点时，新填充边的两个端点仍在\(B_\ell\)中，故继续被该袋覆盖；
从袋中删去已消顶点也不破坏其余顶点的运行交集。因此不变量在消元过程中保持，且每个
被消顶点的尚存邻居数至多为\(|B_\ell|-1\leq k\)。当\(X\)删空后，叶袋只剩接口
\(B_\ell\cap B_p\)，其中的边也由父袋\(B_p\)覆盖，因而可以移除该叶袋并保持不变量。
重复到只剩根袋，再以任意次序消去根袋。由此，树宽等于所有消元次序的最小诱导宽度。
寻找最优次序本身通常很难，实践中常用最小度数或最少填充边等启发式，并把所得宽度作为
可计算上界。

\section{分配律与动态规划的代数}
\label{sec:graph-theory-semiring-dynamic-programming}

\subsection{共享子问题依靠分配律避免重复枚举}

考虑贯穿图上的表达式~\eqref{eq:graph-theory-factorized-objective}。若要汇总全部赋值，
直接写成
\[
  Z=\bigoplus_{z_1,\ldots,z_4}F(\bm{z}).
\]
变量\(z_1\)只通过接口\((z_2,z_3)\)与其余变量联系，因此可以先计算
\begin{equation}
  m_{1\to\{2,3\}}(z_2,z_3)=
  \bigoplus_{z_1\in\mathcal{Z}_1}
  \phi_1(z_1)\otimes
  \phi_{1,2}(z_1,z_2)\otimes\phi_{1,3}(z_1,z_3).
  \label{eq:graph-theory-interface-message}
\end{equation}
对每个接口赋值\((z_2,z_3)\)，这个局部结果只需计算一次。变量\(z_4\)一侧可得到
同形消息\(m_{4\to\{2,3\}}(z_2,z_3)\)，最终只需在接口上组合两条消息和其余局部项。

这一步合法，不是因为图画成树枝形，而是因为\(\otimes\)对\(\oplus\)满足分配律。
若把与\(z_1\)无关、但可以依赖接口值的项记为\(c(z_2,z_3)\)，则
\[
  \bigoplus_{z_1}\bigl(c(z_2,z_3)\otimes g(z_1,z_2,z_3)\bigr)
  =c(z_2,z_3)\otimes\left(\bigoplus_{z_1}g(z_1,z_2,z_3)\right).
\]
图说明哪些项与\(z_1\)无关，代数规则则允许把它们移出汇总。动态规划的核心就是保存
接口变量上的局部结果，使后续不同部分复用同一个子问题答案。

树上的消息可以把这一直觉写成精确不变量。先考虑一棵树，并在顶点与边上分别放置
\(\phi_i(z_i)\)和\(\phi_{i,j}(z_i,z_j)\)。删除边\(\{u,v\}\)后，记包含\(u\)
的连通分量为\(T_{u\mid v}\)。从\(u\)发给\(v\)的消息定义为
\begin{equation}
  m_{u\to v}(z_v)=
  \bigoplus_{\bm{z}_{T_{u\mid v}}}
  \left[
    \bigotimes_{i\in T_{u\mid v}}\phi_i(z_i)
    \otimes
    \bigotimes_{\substack{\{i,j\}\in E\\i,j\in T_{u\mid v}}}
      \phi_{i,j}(z_i,z_j)
    \otimes\phi_{u,v}(z_u,z_v)
  \right].
  \label{eq:graph-theory-tree-message-semantics}
\end{equation}
这里消息的自变量只有接口值\(z_v\)，被汇总的范围则是整棵\(u\)侧子树。

\begin{theorem}[树消息递推]
  \label{thm:graph-theory-tree-message-recursion}
  在有限树和交换半环上，式~\eqref{eq:graph-theory-tree-message-semantics}满足
  \begin{equation}
    m_{u\to v}(z_v)=
    \bigoplus_{z_u}
    \left[
      \phi_u(z_u)\otimes\phi_{u,v}(z_u,z_v)
      \otimes
      \bigotimes_{r\in N(u)\setminus\{v\}}m_{r\to u}(z_u)
    \right].
    \label{eq:graph-theory-tree-message-recursion}
  \end{equation}
\end{theorem}

\begin{proof}
对子树\(T_{u\mid v}\)的顶点数归纳。若\(u\)是叶顶点，乘积中没有入消息，递推式正是
对\(z_u\)汇总顶点项与接口边项。若\(u\)还有邻居\(r\neq v\)，删除\(u\)后各
\(T_{r\mid u}\)互不相交，且它们之间没有边。归纳假设把每棵子树的全部局部项概括为
\(m_{r\to u}(z_u)\)。结合律和交换律把这些子树分组，分配律再把只依赖\(z_u,z_v\)
的项移到各子树汇总之外，便得到式
~\eqref{eq:graph-theory-tree-message-recursion}。
\end{proof}

定理依赖树在删边后确实分成互不重叠的子树。一般有环图若直接沿原边反复发送同形消息，
同一局部项可能经不同方向重复返回；要得到精确结果，须先消元或构造满足运行交集的树分解。

\subsection{半环抽取局部计算真正需要的规则}

\term{半环}（semiring）
\((\mathcal{S},\oplus,\otimes,e_{\oplus},e_{\otimes})\)由一个集合、两种运算和两个
单位元组成。本章使用交换半环，它满足：
\begin{enumerate}
  \item \(\oplus\)结合、交换，单位元为\(e_{\oplus}\)；
  \item \(\otimes\)结合、交换，单位元为\(e_{\otimes}\)；
  \item \(\otimes\)对\(\oplus\)分配；
  \item \(e_{\oplus}\)对\(\otimes\)具有吸收性，
    \(a\otimes e_{\oplus}=e_{\oplus}\)。
\end{enumerate}
半环不要求每个元素有加法逆元或乘法逆元。动态规划通常只需要组合、汇总和分配，不需要
做减法或除法。

三种与本章直接相关的半环为：
\[
\begin{array}{c|c|c|c|c}
  \text{用途}&\mathcal{S}&\oplus&\otimes&(e_{\oplus},e_{\otimes})\\
  \hline
  \text{总权重}&\mathbb{R}_{\geq0}&+&\times&(0,1)\\
  \text{最大权重}&\mathbb{R}_{\geq0}&\max&\times&(0,1)\\
  \text{最小成本}&\mathbb{R}\cup\{+\infty\}&\min&+&(+\infty,0)
\end{array}
\]
最大积中\(0\)既是\(\max\)的单位元，也是乘法吸收元。最小加法中
\(+\infty+c=+\infty\)，所以\(+\infty\)表示不可行路径或赋值。若最大积取对数，
还可得到\((\max,+)\)半环，其汇总单位元为\(-\infty\)、组合单位元为\(0\)。

布尔半环\((\{\mathrm{false},\mathrm{true}\},\lor,\land)\)则计算是否存在可行赋值。
同一个图结构因此能够回答“总共有多少权重”“最佳配置是什么”“是否存在配置”等不同
查询。结构相同不表示答案相同；运算是问题定义的一部分。

一个两变量小表可以直接核对运算改变了什么。令
\[
  g(x,y)=
  \begin{array}{c|cc}
    &y=0&y=1\\ \hline
    x=0&1&3\\
    x=1&2&4
  \end{array}.
\]
把四个表项视为非负权重时，和积查询给出
\(\sum_{x,y}g(x,y)=10\)，最大积查询在这里只有一个局部因子，给出
\(\max_{x,y}g(x,y)=4\)及见证\((1,1)\)。把同一组数字改解释为成本，最小加法查询
给出\(\min_{x,y}g(x,y)=1\)及见证\((0,0)\)。这三个答案来自同一作用域，却对应
不同问题；数值表本身不能决定应选哪一种半环。

\subsection{变量消除在任意交换半环上复用同一结构}

设每个局部因子取值于半环\(\mathcal{S}\)。给定消元次序，每次收集所有含变量\(z_i\)
的因子，先用\(\otimes\)组合，再对\(z_i\)用\(\oplus\)汇总，得到作用于剩余邻居的
新因子。算法重复到不再有变量，就得到全局半环和
\begin{equation}
  Z=\bigoplus_{\bm{z}\in\mathcal{Z}_1\times\cdots\times\mathcal{Z}_4}
  \left[
    \bigotimes_i\phi_i(z_i)
    \otimes
    \bigotimes_{\{i,j\}\in E}\phi_{i,j}(z_i,z_j)
  \right].
  \label{eq:graph-theory-semiring-elimination}
\end{equation}

\begin{theorem}[半环变量消除的正确性]
  \label{thm:graph-theory-semiring-elimination}
  在有限变量、交换半环和式
  ~\eqref{eq:graph-theory-factorized-objective}的局部分解下，按任意完整次序执行
  上述变量消除，所得结果等于式
  ~\eqref{eq:graph-theory-semiring-elimination}的直接枚举值。
\end{theorem}

\begin{proof}
证明维持一个不变量：每次消元后，当前因子的半环乘积再对剩余变量汇总，等于原表达式
对已消与未消变量的全局汇总。初始时该命题就是式
~\eqref{eq:graph-theory-semiring-elimination}。

设下一变量为\(z_i\)，把当前因子分成含\(z_i\)的集合\(\mathcal{F}_i\)和不含它的
集合\(\mathcal{F}_{-i}\)。有限性、结合律与交换律允许重排有限汇总，分配律给出
\begin{align*}
  &\bigoplus_{\bm{z}_{-i}}\bigoplus_{z_i}\left[
    \left(\bigotimes_{f\in\mathcal{F}_{-i}}f\right)
    \otimes
    \left(\bigotimes_{f\in\mathcal{F}_i}f\right)
  \right]\\
  &\quad=\bigoplus_{\bm{z}_{-i}}\left[
    \left(\bigotimes_{f\in\mathcal{F}_{-i}}f\right)
    \otimes
    \left(
      \bigoplus_{z_i}
      \bigotimes_{f\in\mathcal{F}_i}f
    \right)
  \right].
\end{align*}
括号中的最后一项正是算法生成的新因子，所以消去\(z_i\)后不变量保持。归纳到所有变量
消去，剩余标量等于直接枚举值。
\end{proof}

定理说明消元次序不改变精确答案，却会改变中间因子的作用域。若所有变量取值数不超过
\(q\)，某次序诱导宽度为\(k\)，稠密表实现的时间通常含
\(O(|V|q^{k+1})\)量级，空间含\(O(|V|q^k)\)量级；具体常数还取决于因子数和存储
方式。树宽小使指数只依赖局部宽度，不代表变量取值很大或因子计算昂贵时成本一定低。

\subsection{和积、最大积与最小加法改变所求答案}

若\(\phi\)是非负权重，选择和积半环得到总权重\(Z\)。固定某个变量值而汇总其余变量，
可以得到未归一化边缘权重；若这些权重来自概率模型，还需除以总权重并核对分布定义。
半环公式本身只计算非负数，不自动赋予概率解释。

改用最大积半环，
\[
  M=\max_{\bm{z}}\prod_i\phi_i(z_i)
  \prod_{\{i,j\}\in E}\phi_{i,j}(z_i,z_j)
\]
得到最佳权重。若还要恢复最优配置，每次局部最大化必须保存达到最大值的\(z_i\)选择，
最后按消元逆序回溯。只保存最大值会丢掉见证赋值；并列最大值时，还应说明返回任一个、
全部，还是采用固定平局规则。

若局部量是成本\(c_i,c_{i,j}\)，采用最小加法半环得到
\[
  C^*=\min_{\bm{z}}\left[
    \sum_i c_i(z_i)+\sum_{\{i,j\}\in E}c_{i,j}(z_i,z_j)
  \right].
\]
这与最大积取负对数的形式相近，但只有权重为正并且对数变换保持目标次序时才能对应。
零权重映为无穷成本，负权重则不能直接取实对数。

最短路需要先说明在哪一类对象上优化。记\(\mathcal{P}_{s,v}\)为从\(s\)到\(v\)
的简单路径集合，\(\mathcal{W}_{s,v}\)为允许重复顶点和边的有限游走集合，并约定空集
上的最小值或下确界为\(+\infty\)。对应的两个量为
\[
  d_{\mathrm{path}}(s,v)
  =\min_{P\in\mathcal{P}_{s,v}}\ell(P),
  \qquad
  d_{\mathrm{walk}}(s,v)
  =\inf_{W\in\mathcal{W}_{s,v}}\ell(W).
\]
有限图中的简单路径只有有限条，所以目标可达时\(d_{\mathrm{path}}\)总能由某条路径
达到。游走可以反复绕环，因而\(d_{\mathrm{walk}}\)可能为\(-\infty\)。若不存在
从\(s\)可达且还能到达\(v\)的负环，则可以从任意游走中删除非负闭合段而不增加长度，
最终得到简单路径，所以两个量相等；若存在这样的负环，则
\(d_{\mathrm{walk}}(s,v)=-\infty\)，而简单路径最小值仍是有限条候选中的最小值。

标准最短路算法通常求游走下确界。先假设不存在从\(s\)可达的负环，此时每个有限下确界
都可由简单路径见证。把从起点\(s\)到顶点\(v\)的距离记为\(D(v)\)。空游走给出
起点边界\(D(s)=0\)；若\(v\)从\(s\)不可达，则约定\(D(v)=+\infty\)。对其余顶点，
每条非空路径都可按最后一条边\(u\to v\)分组，于是得到Bellman最优性关系
\begin{equation}
  D(s)=0,\qquad
  D(v)=\min_{u\to v}\{D(u)+\ell_{u,v}\}\quad(v\neq s),
  \label{eq:graph-theory-shortest-path-bellman}
\end{equation}
其中空前驱集合上的最小值为\(+\infty\)，并约定
\(+\infty+c=+\infty\)。这些约定使不可达顶点无需另写例外分支。

在DAG上，式~\eqref{eq:graph-theory-shortest-path-bellman}确实给出一次递推。先令
\(D(s)=0\)、其余顶点为\(+\infty\)，再按拓扑序处理顶点\(v\)，用所有入边计算
\[
  D(v)\leftarrow
  \min\left\{D(v),\ \min_{u\to v}\bigl(D(u)+\ell_{u,v}\bigr)\right\}.
\]
当处理\(v\)时，它的全部前驱已经处理。任意到达\(v\)的非空路径都以某条
\(u\to v\)结束，而每个候选\(D(u)+\ell_{u,v}\)也对应一条实际路径；对拓扑序归纳
便知更新后的\(D(v)\)恰为最短距离。不可达顶点的所有候选仍为\(+\infty\)，因此保持
不可达。每条边只参与一次松弛，路径与游走也因无环而没有区别。

一般有环图不能把Bellman式按任意顶点次序只计算一遍，因为一个前驱的值可能稍后继续
下降。甚至在没有负环时，循环固定点方程本身也可能有多组解。例如加入从\(s\)到
\(a\)的长度\(5\)边以及\(a,b\)之间长度均为\(0\)的双向边后，任意
\(D(a)=D(b)\leq5\)都满足局部等式，但真正的最短距离是\(5\)。算法还必须规定从
\(+\infty\)开始的松弛过程，而不能只写固定点方程。

为处理尚不知道是否存在负环的一般输入，改用迭代标签而不预先称其为真实距离。令
\(L^{(0)}(s)=0\)，对\(v\neq s\)令\(L^{(0)}(v)=+\infty\)，并递推
\[
  L^{(k)}(v)=
  \min\left\{
    L^{(k-1)}(v),\
    \min_{u\to v}\bigl(L^{(k-1)}(u)+\ell_{u,v}\bigr)
  \right\}.
\]
对\(k\)归纳可知，\(L^{(k)}(v)\)是使用至多\(k\)条边到达\(v\)的最小游走长度。
若不存在从\(s\)可达的负环，则每个有限最短距离都有至多\(|V|-1\)条边的简单路径
见证，故\(|V|-1\)轮足够；这正是Bellman--Ford松弛的基础。若某个从\(s\)可达的
负环还能到达目标\(v\)，则\(L^{(k)}(v)\)会因反复绕环而无下界，真实游走下确界为
\(-\infty\)，不存在有限的Bellman距离。下一节将分别说明非负边长怎样支持Dijkstra
的贪心次序，以及怎样先用额外一轮松弛找到可达负环的证据，再沿有向边传播其影响。

\subsection{求和与取最大混用时不能任意换序}

同一种\(\oplus\)在有限变量间可以交换次序，但不同汇总运算通常不能。取
\[
  h(x,y)=
  \begin{array}{c|cc}
    &y=0&y=1\\ \hline
    x=0&4&0\\
    x=1&0&4
  \end{array}.
\]
若先对每个\(x\)把\(y\)求和，再选同一个\(x\)，则
\begin{equation}
  \max_x\sum_y h(x,y)=4.
  \label{eq:graph-theory-max-sum-order-a}
\end{equation}
若允许每个\(y\)分别选择最好的\(x\)，则
\begin{equation}
  \sum_y\max_x h(x,y)=8.
  \label{eq:graph-theory-max-sum-order-b}
\end{equation}
两个值不同，因为第二个查询允许选择随\(y\)变化，而第一个查询要求在看到或汇总\(y\)
之前固定同一个\(x\)。

一般总有
\[
  \max_x\sum_y h(x,y)\leq\sum_y\max_x h(x,y),
\]
因为对任意固定\(x\)，每一项\(h(x,y)\)都不超过\(\max_{x'}h(x',y)\)；再对\(x\)
取最大仍保持不等式。等号需要存在一个共同的\(x\)同时在所有相关\(y\)上达到最大，
不能只凭分配律假定。

因此，含求和变量与最大化变量的查询不仅由图和树宽决定，还由允许的消元次序决定。
受约束次序可能比普通树宽宽得多。把边缘概率、最可能联合配置和边缘最大化混称为一次
“推断”，会掩盖它们不同的查询目标和计算边界。

前述交换与分配论证针对有限汇总。把\(\oplus\)换成连续积分时，还需第
~\ref{sec:analysis-limit-integral-interchange}节的Tonelli或Fubini条件：非负可测
函数可用Tonelli交换积分，带符号函数通常需要绝对可积。仅仅把有限和积公式中的求和号
改成积分号，不能保证积分存在或允许换序。类似地，树消息的精确性来自无环分解；在原始
环图上运行局部传播可以成为近似算法，但其收敛与误差不由半环公理自动保证。

\section{图上的精确与近似求解}
\label{sec:graph-theory-exact-approximate-algorithms}

\subsection{最短路算法依赖边长条件}

给有向或无向图的每条边指定实边长\(\ell_e\)。本节的单源距离采用上文的
\(d_{\mathrm{walk}}(s,v)\)；距离有限时，它等于简单路径最小值。若所有边长非负，
Dijkstra算法维护已确定集合\(S\)和暂定距离\(D(v)\)：每次选择\(S\)外暂定距离
最小的顶点\(u\)，把它加入\(S\)，再尝试用\(D(u)+\ell_{u,v}\)改善邻居
\scite{math-cormen2022algorithms}

\begin{theorem}[Dijkstra算法的贪心正确性]
  \label{thm:graph-theory-dijkstra-correctness}
  若所有边长非负，则顶点\(u\)被选入已确定集合时，暂定距离\(D(u)\)等于从\(s\)
  到\(u\)的游走下确界，也等于简单路径最小值。
\end{theorem}

\begin{proof}
对选入次序归纳。起点距离为零，结论成立。假设此前选入顶点均已正确，考虑新选入的
\(u\)。若存在更短路径\(P\)，沿\(P\)从\(s\)前进，令\(y\)为第一个尚未进入\(S\)
的顶点，\(x\)为它在路径上的前驱。归纳假设给出\(D(x)\)已经等于最短距离；处理
\(x\)时，松弛保证
\[
  D(y)\leq D(x)+\ell_{x,y},
\]
不超过路径\(P\)到达\(y\)的前缀长度。由于剩余边长非负，该前缀长度不超过\(P\)
到\(u\)的总长度。算法又因\(u\)的暂定距离最小而有\(D(u)\leq D(y)\)。于是
\(D(u)\)不超过假设中的更短路径长度，与\(D(u)\)始终由某条实际路径产生、不能低于
最短距离矛盾。
\end{proof}

非负性用在“路径前缀不会比总长度更大”这一步。存在负边时，一个已确定顶点可能在后面
通过负边得到更低成本，贪心结论失效。Bellman--Ford算法反复松弛全部边；经过
\(|V|-1\)轮后，它能得到所有不受相关负环影响的有限距离。接着检查一轮所有边：若
\(L(u)<+\infty\)且\(L(u)+\ell_{u,v}<L(v)\)，就把可继续改善的端点\(v\)记为
受影响种子。这一轮只检测从\(s\)可达的负环证据，尚未标出全部受影响目标；还要从所有
种子出发沿有向边做一次可达性遍历，所有遍历到的顶点的游走下确界均为\(-\infty\)。
对特定目标\(v\)，只有某个负环既能从\(s\)到达、又能沿有向路径到达\(v\)时，
\(d_{\mathrm{walk}}(s,v)\)才等于\(-\infty\)。若负环从\(s\)可达却不能到达\(v\)，
它不改变\(v\)的距离；若负环本身不能从\(s\)到达，它不影响任何单源查询。检测到负环、
目标距离无有限值与简单路径最小值仍存在，是三个不同判断。

例如取弧长
\[
  \ell_{s,a}=2,\quad \ell_{s,b}=5,\quad
  \ell_{a,b}=1,\quad \ell_{a,t}=5,\quad \ell_{b,t}=2.
\]
Dijkstra算法先确定\(D(a)=2\)，经\(a\)把\(D(b)\)从\(5\)改为\(3\)；随后确定
\(b\)，再把\(D(t)\)从\(7\)改为\(5\)。最短路为\(s\to a\to b\to t\)，长度
\(2+1+2=5\)。这个手算展示的是非负边长下的确定次序；若把\(\ell_{a,b}\)改成负数，
算出的某条路径仍可验证，却不能继续用定理
~\ref{thm:graph-theory-dijkstra-correctness}保证贪心确定步骤。

\subsection{最小生成树连接全部顶点而不优化点对距离}

在连通无向加权图中，\term{生成树}（spanning tree）包含全部顶点并且是一棵树。
\term{最小生成树}（minimum spanning tree, MST）使所选\(n-1\)条边的总权重最小。
它优化全局连接成本，不要求任意两点在树上的路径仍是原图最短路。

MST贪心算法依据割性质。对任意非平凡顶点划分\((S,V\setminus S)\)，跨越该割的
最轻边称为轻边。
\begin{lemma}[最小生成树的割性质]
  \label{lem:graph-theory-mst-cut-property}
  若边集合\(A\)包含在某棵最小生成树中，且割\((S,V\setminus S)\)不切断\(A\)
  中的任何边，则跨割的一条轻边\(e\)也可以与\(A\)共同包含在某棵最小生成树中。
\end{lemma}

\begin{proof}
取一棵包含\(A\)的最小生成树\(T\)。若\(e\in T\)，结论已成立。否则把\(e\)加入
\(T\)会形成唯一环。这个环还必须通过另一条跨割边\(f\)，否则无法从割的一侧经\(e\)
返回同侧。由于\(e\)是轻边，\(w_e\leq w_f\)。删除\(f\)得到新生成树
\(T'=T+e-f\)，其总权重不大于\(T\)，故仍为最小生成树。割不切断\(A\)中的边，
所以\(f\notin A\)，新树仍包含\(A\)。
\end{proof}

Kruskal算法按边权从小到大扫描，只加入连接不同当前分量的边；Prim算法则从一个已连通
顶点集出发，反复加入跨当前割的轻边。二者都可逐步应用引理
\ref{lem:graph-theory-mst-cut-property}，但维护的割不同。边权并列时MST可能不唯一，
最小总权重仍唯一确定。

\subsection{最大流与最小割给出容量证书}

在带源点\(s\)、汇点\(t\)的有向网络中，每条弧有非负容量\(c_{i,j}\)。
\term{流}（flow）\(f_{i,j}\)满足
\[
  0\leq f_{i,j}\leq c_{i,j},
\]
并且除\(s,t\)外每个顶点流入等于流出。流值\(|f|\)是源点净流出量。
一个\(s\)--\(t\)割由包含\(s\)但不含\(t\)的集合\(S\)确定，容量为
\[
  c(S,V\setminus S)=
  \sum_{\substack{i\in S\\j\notin S}}c_{i,j}.
\]
任意流值都不超过任意割容量，因为对\(S\)内顶点的流守恒式求和，内部流相消，跨割
净流出等于\(|f|\)，而正向流出不超过正向容量、反向流入非负。

\begin{theorem}[最大流最小割定理]
  \label{thm:graph-theory-max-flow-min-cut}
  在有限非负容量网络中，最大\(s\)--\(t\)流值等于最小\(s\)--\(t\)割容量。
\end{theorem}

\begin{proof}
可行流集合由有限个闭有界线性约束定义，非空且紧，因此连续的流值函数能达到最大值。
取一个最大流\(f\)。构造残量网络：弧\(i\to j\)有剩余容量
\(c_{i,j}-f_{i,j}\)，并为已发送的流加入容量\(f_{i,j}\)的反向弧。
若残量网络中存在从\(s\)到\(t\)的路径\(P\)，令\(\delta>0\)为路径上最小剩余
容量。沿\(P\)更新原流：经过原弧\(i\to j\)对应的正向残量弧时令
\(f_{i,j}\leftarrow f_{i,j}+\delta\)，经过该原弧对应的反向残量弧时令
\(f_{i,j}\leftarrow f_{i,j}-\delta\)。瓶颈值的定义保证正向更新不超过容量、反向
更新不低于零；路径内部每个顶点的一次流入变化与一次流出变化相消，所以流量守恒保持，
而源点净流出量增加\(\delta\)。这便得到更大可行流，与最大性矛盾。

令\(S\)为残量网络中从\(s\)可达的顶点集合。于是\(s\in S,t\notin S\)。
对每条从\(S\)到补集的原弧，剩余容量必须为零，否则终点也可达，所以
\(f_{i,j}=c_{i,j}\)。对每条从补集进入\(S\)的原弧，其反向残量容量必须为零，所以
\(f_{i,j}=0\)。因此跨割净流出恰等于
\[
  |f|=\sum_{\substack{i\in S\\j\notin S}}c_{i,j}
  =c(S,V\setminus S).
\]
前面的弱上界说明没有流能超过任何割；现在一个最大流与一个割达到同一值，二者分别最优。
\end{proof}

残量网络同时给出最优性证书。对整数容量，沿整数瓶颈增广可得到整数最大流。对有理容量
或机器有限精度容量，容量具有有限比特编码，可以采用具有有限步复杂度保证的最大流算法，
并把位运算成本计入输入长度。对任意数学实数，定理仍由紧性论证保证最优流和最小割存在；
若没有给出这些实数的有限有效表示或精确运算接口，就不能进一步声称程序能够精确输出
最优解。

一个整数网络可以同时展示增广与证书。令
\[
  c_{s,a}=3,\quad c_{s,b}=2,\quad c_{a,b}=1,\quad
  c_{a,t}=2,\quad c_{b,t}=3.
\]
依次沿\(s\to a\to t\)、\(s\to b\to t\)和
\(s\to a\to b\to t\)增广\(2,2,1\)单位，得到流值\(5\)。此时割
\(S=\{s\}\)的容量为\(c_{s,a}+c_{s,b}=5\)，所以弱上界已经达到：无需枚举其他流，
便可由“可行流值\(5\)+割容量\(5\)”认证二者分别为最大流和最小割。

\subsection{二部匹配是流的一个精确特例}

二部图\(\mathcal{G}=(U,W,E)\)中的\term{匹配}（matching）是一组互不共享端点的边。
要求边数最大的匹配称为最大匹配。建立源点到每个\(u\in U\)容量为\(1\)的弧，把每条
\(\{u,w\}\in E\)改为容量\(1\)的弧\(u\to w\)，再从每个\(w\in W\)连容量\(1\)
的弧到汇点。任一整数流值为\(k\)对应\(k\)条匹配边，反之亦然。整数容量的最大流存在
整数最优解，所以二部最大匹配可由最大流精确求得。

匹配语言中的\term{增广路径}（augmenting path）是一条端点尚未匹配、边在“未匹配”
与“已匹配”之间交替出现的路径。沿它翻转两类边会使匹配大小增加\(1\)。在上述流网络
中，这正对应一条残量\(s\)--\(t\)路径；不存在增广路径时，残量割给出最大性的证书。
这个接口只说明基数二部匹配与单位容量流的对应，不涵盖一般图匹配所需的奇环处理。

非二部图的奇环会使这个简单流构造失效，加权匹配又改变目标。“可画成图”不足以选择
算法，还要核对图类别、约束与优化目标。

\subsection{树宽小使局部动态规划精确可行}

树分解上的动态规划为每个袋维护接口赋值的最优值或总权重。以贯穿图为例，叶袋
\(\{1,2,3\}\)向袋\(\{2,3,4\}\)发送关于接口\((z_2,z_3)\)的表；后一袋组合消息
与含\(z_4\)的局部项，再完成余下汇总。运行交集条件保证某变量在不同袋中的出现沿树
保持一致。

若树宽为\(k\)、每个变量至多有\(q\)个取值，袋表规模至多为\(q^{k+1}\)。
这与第~\ref{sec:graph-theory-traversal-separation-decomposition}节的诱导宽度分析是
同一局部表成本，此处不再重复按全部变量数枚举候选。固定\(k\)时，动态规划关于顶点数
为多项式时间；\(k\)随顶点数增长时，这一保证不再成立。

\subsection{二值成对次模能量可以精确转成最小割}

现在回到第~\ref{sec:comb-optimization-submodularity}节留下的二值离散目标。设
\(z_i\in\{0,1\}\)，能量为
\begin{equation}
  \mathcal{E}(\bm{z})=\sum_i D_i(z_i)+
  \sum_{\{i,j\}\in E}\theta_{i,j}(z_i,z_j).
  \label{eq:graph-theory-binary-pairwise-energy}
\end{equation}
这里用\(\mathcal{E}\)表示能量函数，\(E\)仍专指图的边集合。
成对项若满足
\begin{equation}
  \theta_{i,j}(0,0)+\theta_{i,j}(1,1)
  \leq
  \theta_{i,j}(0,1)+\theta_{i,j}(1,0),
  \label{eq:graph-theory-pairwise-submodularity}
\end{equation}
就称该二值成对项为次模。它限制四个配置的相对组合，并不要求相同标签的两个能量分别
都小于不同标签能量。

\begin{theorem}[二值次模成对能量的割表示]
  \label{thm:graph-theory-submodular-cut-representation}
  若式~\eqref{eq:graph-theory-binary-pairwise-energy}中所有一元项和成对项均为有限
  实数，且每个成对项满足式
  ~\eqref{eq:graph-theory-pairwise-submodularity}，则存在非负容量的
  \(s\)--\(t\)网络，使每个二值配置对应一个割，割容量等于该配置能量加同一个常数。
  因而一次最小割给出全局最优配置。
\end{theorem}

\begin{proof}
固定一条边并简写
\(\theta_{a,b}=\theta_{i,j}(a,b)\)。令
\begin{align}
  w_{i,j}&=\frac{
    \theta_{0,1}+\theta_{1,0}
    -\theta_{0,0}-\theta_{1,1}
  }{2},
  \label{eq:graph-theory-cut-pair-weight}\\
  \alpha_i&=\theta_{1,0}-\theta_{0,0}-w_{i,j},
  \qquad \alpha_j=
  \theta_{0,1}-\theta_{0,0}-w_{i,j}.
  \label{eq:graph-theory-cut-unary-shifts}
\end{align}
次模不等式恰好保证\(w_{i,j}\geq0\)。逐一代入四种二值配置可验证
\begin{equation}
  \theta_{i,j}(z_i,z_j)=
  \theta_{0,0}+\alpha_i z_i+\alpha_j z_j+
  w_{i,j}|z_i-z_j|.
  \label{eq:graph-theory-pairwise-cut-decomposition}
\end{equation}
具体地，
\begin{align*}
  (0,0):\quad&
  \theta_{0,0}=\theta_{0,0},\\
  (1,0):\quad&
  \theta_{0,0}+\alpha_i+w_{i,j}=\theta_{1,0},\\
  (0,1):\quad&
  \theta_{0,0}+\alpha_j+w_{i,j}=\theta_{0,1},\\
  (1,1):\quad&
  \theta_{0,0}+\alpha_i+\alpha_j
  =\theta_{1,1},
\end{align*}
其中最后一个等式使用
\(2w_{i,j}=\theta_{0,1}+\theta_{1,0}-\theta_{0,0}-\theta_{1,1}\)。
式中的\(\theta_{0,0}\)可并入总常数，线性项并入相应一元函数。对所有边重复后，
能量写成
\[
  C+\sum_i D'_i(z_i)
  +\sum_{\{i,j\}\in E}w_{i,j}|z_i-z_j|.
\]
一元值可能为负，但对每个\(i\)同时减去
\(\min\{D'_i(0),D'_i(1)\}\)只改变常数，并使两个值都非负。

建立源点\(s\)、汇点\(t\)和每个变量顶点。规定源侧表示\(z_i=0\)，汇侧表示
\(z_i=1\)。加入容量\(D'_i(1)\)的弧\(s\to i\)和容量\(D'_i(0)\)的弧
\(i\to t\)。若\(i\)在源侧，只有\(i\to t\)跨割，支付\(D'_i(0)\)；若在汇侧，
只有\(s\to i\)跨割，支付\(D'_i(1)\)。对每个无向成对项
\(\{i,j\}\)，加入两条容量均为\(w_{i,j}\)的反向弧\(i\to j\)与\(j\to i\)。
当两端同侧时没有弧跨割；当两端异侧时，恰有一条弧从源侧指向汇侧，因而无论
\((z_i,z_j)=(0,1)\)还是\((1,0)\)，割容量都增加
\(w_{i,j}|z_i-z_j|=w_{i,j}\)。

因此每个配置的割容量等于\(\mathcal{E}(\bm{z})-C'\)，其中\(C'\)与配置无关。每个
\(s\)--\(t\)割也唯一规定所有变量位于哪一侧，所以最小割与全局最小能量配置一一对应。
\end{proof}

这一结论对有限实值项给出割表示和最优解的存在性。若各项是有理数或机器有限精度数，
构造出的容量具有有限编码，因而可调用最大流算法计算最优配置；对没有有效表示的任意
实数，只保留存在性结论。二元成对能量可由图割表示的更一般条件见Kolmogorov和Zabih
\scite{pgm-variational-kolmogorov2004}

取两个变量，令
\[
  D_1(0)=0,\quad D_1(1)=2,\qquad
  D_2(0)=2,\quad D_2(1)=0,\qquad
  w_{1,2}=1.
\]
四个配置\((0,0),(0,1),(1,0),(1,1)\)的能量依次为
\(2,1,5,2\)。最小割把顶点\(1\)放在源侧、顶点\(2\)放在汇侧，只支付二者间容量
\(1\)，返回最优配置\((0,1)\)。这个小例子同时核对了终端边的方向。

若式~\eqref{eq:graph-theory-pairwise-submodularity}失败，
\(w_{i,j}<0\)，非负容量割网络不能按上述方式表示该项。多标签能量、含高阶项的能量或
一般非次模二次目标也不自动归入一次最小割。某些特殊结构可通过辅助顶点、分层网络或
一系列二值子问题处理，但必须分别证明等价性或近似比。

\subsection{松弛、取整与搜索需要报告界}

精确结构不存在时，可以把二值约束\(z_i\in\{0,1\}\)放宽为
\(z_i\in[0,1]\)，但扩大定义域还不足以定义松弛问题。记原离散可行集为
\(\mathcal{Z}\)，还需同时给出包含它的连续可行域
\(\widetilde{\mathcal{Z}}\supseteq\mathcal{Z}\)，以及连续域上的目标延拓
\(\widetilde{\mathcal{E}}\)，并要求
\[
  \widetilde{\mathcal{E}}(\bm{z})=\mathcal{E}(\bm{z})
  \qquad(\bm{z}\in\mathcal{Z}).
\]
固定这样的延拓后，
\[
  L_{\mathrm{relax}}
  =\inf_{\bm{z}\in\widetilde{\mathcal{Z}}}
    \widetilde{\mathcal{E}}(\bm{z})
  \leq
  \min_{\bm{z}\in\mathcal{Z}}\mathcal{E}(\bm{z}),
\]
所以松弛下确界是原最小化问题的下界。若松弛问题达到该下确界，且一个松弛最优解恰好
属于\(\mathcal{Z}\)，它同时是原问题的全局最优解；若为分数解，还需取整、局部搜索
或分支定界。不同目标延拓即使在全部二值点上一致，也可能给出不同的松弛下界，报告结果
时必须说明所用延拓。

取整得到的是一个原问题可行解及其目标上界，松弛值给出下界。二者差距
\[
  \mathcal{E}(\widehat{\bm{z}})-L_{\mathrm{relax}}
\]
是可计算的次优证书，但不等于真实误差，除非下界已经紧。逐点按\(1/2\)阈值取整可能
破坏全局约束；随机取整的期望保证也必须说明随机性、可行性修复和适用目标。

近似算法应给出比较口径。乘法近似比适合非负且最优值不为零的目标；目标可为负或跨过
零时，加法误差更合适。实例上运行很快不构成复杂度或近似保证，图结构、目标与证书必须
一起报告。

\section{图谱与离散对称性}
\label{sec:graph-theory-spectrum-discrete-symmetry}

\subsection{图拉普拉斯把边上的差分汇总到顶点}

考虑无向非负加权图，并约定非边的权重为\(0\)。沿用前文记号，
\(E_+=\{\{i,j\}\in E:w_{i,j}>0\}\)，支撑图为
\(\mathcal{G}_+=(V,E_+)\)；本节所说的支撑连通分量均指这张图的连通分量。令
\(\bm{W}=[w_{i,j}]\)为对称权重矩阵，
\(d_i^{(w)}=\sum_jw_{i,j}\)为顶点\(i\)的\term{加权度}（weighted degree），并令
\(\bm{D}^{(w)}=\operatorname{diag}(d_1^{(w)},\ldots,d_n^{(w)})\)。
它不同于第~\ref{sec:graph-theory-objects-representations}节定义的普通度
\(d_i=|N(i)|\)；只有无权图取\(w_{i,j}=1\)时二者才相等
\scite{clust-chung1997}
\term{组合图拉普拉斯}（combinatorial graph Laplacian）定义为
\begin{equation}
  \bm{L}=\bm{D}^{(w)}-\bm{W}.
  \label{eq:graph-theory-combinatorial-laplacian}
\end{equation}
对顶点信号\(\bm{f}\in\mathbb{R}^n\)，
\[
  (\bm{L}\bm{f})_i=\sum_jw_{i,j}(f_i-f_j),
\]
它比较顶点值与邻居值，并按边权汇总差异。若用前文的有向关联矩阵\(\bm{B}\)，把边权
放入对角矩阵\(\bm{W}_E\)，则
\[
  \bm{L}=\bm{B}\bm{W}_E\bm{B}^{\top}.
\]
任取边方向只改变\(\bm{B}\)相应列的符号，乘积不变。

\begin{theorem}[图拉普拉斯的Dirichlet能量与零空间]
  \label{thm:graph-theory-laplacian-energy-nullspace}
  对有限无向非负加权图，
  \begin{equation}
    \bm{f}^{\top}\bm{L}\bm{f}
    =\frac12\sum_{i,j}w_{i,j}(f_i-f_j)^2\geq0.
    \label{eq:graph-theory-dirichlet-energy}
  \end{equation}
  因而\(\bm{L}\)对称半正定。其零空间维数等于正权边构成图的连通分量数，并由各分量的
  指示向量张成。
\end{theorem}

\begin{proof}
展开二次型，
\begin{align*}
  \bm{f}^{\top}\bm{L}\bm{f}
  &=
  \sum_i d_i^{(w)}f_i^2-\sum_{i,j}w_{i,j}f_if_j\\
  &=
  \frac12\sum_{i,j}w_{i,j}
  \bigl(f_i^2+f_j^2-2f_if_j\bigr),
\end{align*}
其中利用了\(w_{i,j}=w_{j,i}\)以及
\(\sum_jw_{i,j}=d_i^{(w)}\)。这就得到式
~\eqref{eq:graph-theory-dirichlet-energy}和半正定性。

若\(\bm{L}\bm{f}=\bm{0}\)，则二次型为零。每一项非负，所以每条正权边都满足
\(f_i=f_j\)。沿正权支撑图中的路径传递可知，同一支撑连通分量内\(\bm{f}\)为常数。
反过来，在每个支撑连通分量内取任意常数时，每条正权边两端值相同，而零权边不贡献
\(\bm{L}\bm{f}\)，故\(\bm{L}\bm{f}=\bm{0}\)。各支撑连通分量的指示向量线性无关，
并生成全部这类分段常数向量，零空间维数即为支撑连通分量数。
\end{proof}

负边权会使半正定结论失效；有向图也有多种不等价的拉普拉斯定义。本节谱结论只针对
无向非负权图。

顶点度差异很大时，常改用\term{对称归一化拉普拉斯}（symmetric normalized graph
Laplacian）
\begin{equation}
  \bm{L}_{\mathrm{sym}}
  =(\bm{D}^{(w)})^{-1/2}\bm{L}(\bm{D}^{(w)})^{-1/2}.
  \label{eq:graph-theory-normalized-laplacian}
\end{equation}
这里对\(d_i^{(w)}=0\)的孤立点约定相应逆平方根为\(0\)，因而孤立点对应的行列也为
零。对正度顶点，令\(\bm{g}=(\bm{D}^{(w)})^{-1/2}\bm{f}\)，便有
\(\bm{f}^{\top}\bm{L}_{\mathrm{sym}}\bm{f}
=\bm{g}^{\top}\bm{L}\bm{g}\geq0\)。正度连通分量上的零特征向量不再是
\(\bm{1}\)，而是\((\bm{D}^{(w)})^{1/2}\bm{1}\)在该分量上的限制。

另一常见形式是随机游走拉普拉斯
\(\bm{L}_{\mathrm{rw}}=(\bm{D}^{(w)})^{-1}\bm{L}\)。它在普通欧氏内积下通常不对称，
但在所有顶点度为正时与\(\bm{L}_{\mathrm{sym}}\)相似，并在加权内积
\(\langle\bm{x},\bm{y}\rangle_{\bm{D}}=\bm{x}^{\top}\bm{D}^{(w)}\bm{y}\)
下自伴。组合、对称归一化与随机游走拉普拉斯对应不同度量；引用谱界或构造扩散前必须
说明采用哪一种，不能混用其特征向量和孤立点约定。

\subsection{低频特征向量在强边上变化较小}

由第~\ref{chap:matrix-analysis}章的谱定理，\(\bm{L}\)有标准正交特征向量
\(\bm{u}_1,\ldots,\bm{u}_n\)，特征值满足
\[
  0=\lambda_1\leq\lambda_2\leq\cdots\leq\lambda_n.
\]
若正权支撑图\(\mathcal{G}_+\)连通，\(\bm{u}_1=\bm{1}/\sqrt{n}\)，且
\(\lambda_2>0\)。
Rayleigh商
\[
  \frac{\bm{f}^{\top}\bm{L}\bm{f}}
       {\bm{f}^{\top}\bm{f}}
\]
衡量单位能量信号沿边变化多少。排除常量方向后使该商最小的向量是
\(\bm{u}_2\)，最小值为\(\lambda_2\)。正权支撑图连通时，这个第二特征向量称为
费德勒向量（Fiedler vector）。它在高权边两端倾向于取相近值，因此可作为图上的
低频坐标。

若取集合\(S\subset V\)，令\(s_i=\mathbb{I}\{i\in S\}\)，并记
\(\bm{s}=(s_1,\ldots,s_n)^{\top}\)。这里与前章一致，\(\mathbb{I}\{\cdot\}\)
表示条件成立时取\(1\)、否则取\(0\)的指示函数。由式
~\eqref{eq:graph-theory-dirichlet-energy}可得
\begin{equation}
  \bm{s}^{\top}\bm{L}\bm{s}
  =\sum_{\substack{i\in S\\j\notin S}}w_{i,j}
  =
  \operatorname{cut}(S,V\setminus S).
  \label{eq:graph-theory-indicator-cut-energy}
\end{equation}
离散分割因此可以写成二次能量最小化。但若不限制集合大小，取\(S=\varnothing\)或
\(S=V\)得到零割，问题没有意义。下面分别推导按顶点数平衡的比率割和按加权度平衡的
归一化割；二者的离散目标相似，松弛后的特征问题却不同。

先看\term{比率割}（RatioCut）。记\(\bar S=V\setminus S\)，并设
\(a=|S|>0\)、\(b=|\bar S|>0\)，于是\(a+b=n\)。比率割定义为
\begin{equation}
  \operatorname{RatioCut}(S,\bar S)
  =
  \operatorname{cut}(S,\bar S)
  \left(\frac1a+\frac1b\right).
  \label{eq:graph-theory-ratio-cut}
\end{equation}
分母惩罚把极少顶点单独切出的做法。为了把它写成二次型，对分割构造两水平平衡向量
\begin{equation}
  f_i=
  \begin{cases}
    \sqrt{b/a},&i\in S,\\
    -\sqrt{a/b},&i\in\bar S.
  \end{cases}
  \label{eq:graph-theory-ratio-cut-balanced-vector}
\end{equation}
两侧的常数不是任意选择。它们恰好使
\[
  \bm{1}^{\top}\bm{f}
  =a\sqrt{\frac ba}-b\sqrt{\frac ab}=0,
  \qquad
  \bm{f}^{\top}\bm{f}
  =a\frac ba+b\frac ab=n.
\]
同侧边的端点取值相同，不贡献Dirichlet能量；跨割边的端点差为
\[
  \sqrt{\frac ba}+\sqrt{\frac ab}
  =\frac{a+b}{\sqrt{ab}}=\frac{n}{\sqrt{ab}}.
\]
因此
\begin{align}
  \bm{f}^{\top}\bm{L}\bm{f}
  &=
  \frac{n^2}{ab}\operatorname{cut}(S,\bar S)\notag\\
  &=
  n\,\operatorname{RatioCut}(S,\bar S),
  \label{eq:graph-theory-ratio-cut-energy}
\end{align}
其中使用了\(1/a+1/b=n/(ab)\)。最小比率割由此等价于在所有形如式
~\eqref{eq:graph-theory-ratio-cut-balanced-vector}的离散两水平向量中最小化
Rayleigh商。

谱松弛去掉“两水平”限制，只保留平衡和尺度约束：
\begin{equation}
  \min_{\substack{\bm{f}\in\mathbb{R}^n\\
                  \bm{f}^{\top}\bm{1}=0,\quad
                  \bm{f}^{\top}\bm{f}=n}}
  \frac1n\bm{f}^{\top}\bm{L}\bm{f}.
  \label{eq:graph-theory-ratio-cut-relaxation}
\end{equation}
把任意可行向量展开为
\(\bm{f}=\sum_{k=2}^{n}\alpha_k\bm{u}_k\)，正交约束已经去掉
\(\bm{u}_1\)分量，于是
\[
  \frac{\bm{f}^{\top}\bm{L}\bm{f}}
       {\bm{f}^{\top}\bm{f}}
  =
  \frac{\sum_{k=2}^{n}\lambda_k\alpha_k^2}
       {\sum_{k=2}^{n}\alpha_k^2}
  \geq\lambda_2.
\]
等号在\(\bm{f}=\sqrt n\,\bm{u}_2\)时达到。因此费德勒向量解出的是
RatioCut的Rayleigh松弛，\(\lambda_2\)是离散最优值的下界；按符号或阈值把它转回
两侧后，仍须用式~\eqref{eq:graph-theory-ratio-cut}评价所得离散割。

\term{归一化割}（Normalized Cut）改用加权体积
\[
  \operatorname{vol}(S)=\sum_{i\in S}d_i^{(w)}
\]
衡量一侧规模。以下先假设所有顶点的加权度为正，并记
\(p=\operatorname{vol}(S)\)、\(q=\operatorname{vol}(\bar S)\)和
\(\operatorname{vol}(V)=p+q\)。定义
\begin{equation}
  \operatorname{NCut}(S,\bar S)
  =
  \operatorname{cut}(S,\bar S)
  \left(\frac1p+\frac1q\right).
  \label{eq:graph-theory-normalized-cut}
\end{equation}
把式~\eqref{eq:graph-theory-ratio-cut-balanced-vector}中的顶点数换成体积，令
\[
  y_i=
  \begin{cases}
    \sqrt{q/p},&i\in S,\\
    -\sqrt{p/q},&i\in\bar S.
  \end{cases}
\]
直接计算得到
\[
  \bm{y}^{\top}\bm{D}^{(w)}\bm{1}=0,\qquad
  \bm{y}^{\top}\bm{D}^{(w)}\bm{y}=\operatorname{vol}(V),
\]
并且
\[
  \bm{y}^{\top}\bm{L}\bm{y}
  =
  \operatorname{vol}(V)\operatorname{NCut}(S,\bar S).
\]
放宽两水平限制后，Normalized Cut对应的不是普通Rayleigh商，而是
\begin{equation}
  \min_{\substack{\bm{y}\neq\bm{0}\\
          \bm{y}^{\top}\bm{D}^{(w)}\bm{1}=0}}
  \frac{\bm{y}^{\top}\bm{L}\bm{y}}
       {\bm{y}^{\top}\bm{D}^{(w)}\bm{y}}.
  \label{eq:graph-theory-normalized-cut-relaxation}
\end{equation}
把分母固定为\(\operatorname{vol}(V)\)，相应拉格朗日函数可写成
\[
  \mathcal{J}(\bm{y},\lambda,\mu)
  =
  \bm{y}^{\top}\bm{L}\bm{y}
  -\lambda\left(
    \bm{y}^{\top}\bm{D}^{(w)}\bm{y}-\operatorname{vol}(V)
  \right)
  -\mu\bm{y}^{\top}\bm{D}^{(w)}\bm{1}.
\]
对\(\bm{y}\)求驻点条件得到
\[
  2\bm{L}\bm{y}
  -2\lambda\bm{D}^{(w)}\bm{y}
  -\mu\bm{D}^{(w)}\bm{1}
  =\bm{0}.
\]
左乘\(\bm{1}^{\top}\)，利用
\(\bm{1}^{\top}\bm{L}=\bm{0}^{\top}\)和
\(\bm{y}^{\top}\bm{D}^{(w)}\bm{1}=0\)，得到
\(-\mu\operatorname{vol}(V)=0\)，故\(\mu=0\)。驻点条件于是化为广义特征问题
\begin{equation}
  \bm{L}\bm{y}=\lambda\bm{D}^{(w)}\bm{y},
  \qquad
  \bm{y}^{\top}\bm{D}^{(w)}\bm{1}=0.
  \label{eq:graph-theory-normalized-cut-generalized-eigenproblem}
\end{equation}
等价地，令\(\bm{x}=(\bm{D}^{(w)})^{1/2}\bm{y}\)，式
~\eqref{eq:graph-theory-normalized-cut-relaxation}变为
\[
  \min_{\substack{\bm{x}\neq\bm{0}\\
        \bm{x}^{\top}(\bm{D}^{(w)})^{1/2}\bm{1}=0}}
  \frac{\bm{x}^{\top}\bm{L}_{\mathrm{sym}}\bm{x}}
       {\bm{x}^{\top}\bm{x}}.
\]
所以应取\(\bm{L}_{\mathrm{sym}}\)的第二特征向量，再变换回\(\bm{y}\)；它一般不是
组合拉普拉斯\(\bm{L}\)的费德勒向量。当度矩阵为标量倍单位阵，也就是正度图为正则图
时，两类特征方向必然重合；对一般不规则图则不能作此假定。若有孤立顶点，体积可能为零，
必须先单独处理这些顶点或明确采用的零度约定，不能直接除以\(p\)或\(q\)。

四顶点路径给出一个可复算的数值例。令
\[
  1\mathbin{-}2\mathbin{-}3\mathbin{-}4
\]
的三条边权均为\(1\)，则
\[
  \bm{L}=
  \begin{bmatrix}
    1&-1&0&0\\
    -1&2&-1&0\\
    0&-1&2&-1\\
    0&0&-1&1
  \end{bmatrix},
  \qquad
  \bm{D}^{(w)}=\operatorname{diag}(1,2,2,1).
\]
取\(S=\{1,2\}\)，只有边\(\{2,3\}\)跨割。此时
\[
  \operatorname{RatioCut}(S,\bar S)=1,\qquad
  \bm{f}=(1,1,-1,-1)^{\top},\qquad
  \frac{\bm{f}^{\top}\bm{L}\bm{f}}{\bm{f}^{\top}\bm{f}}
  =\frac44=1.
\]
该路径的费德勒特征值和一个对应方向为
\[
  \lambda_2=2-\sqrt2\approx0.586,\qquad
  \bm{u}_2\propto
  (1+\sqrt2,\,1,\,-1,\,-1-\sqrt2)^{\top}.
\]
零阈值恢复分割\(\{1,2\}\mid\{3,4\}\)；枚举四顶点的非平凡二分可确认其离散
RatioCut最优值为\(1\)，而松弛下界为\(2-\sqrt2\)。

同一分割两侧体积均为\(3\)，所以
\[
  \operatorname{NCut}(S,\bar S)=\frac23,\qquad
  \frac{\bm{f}^{\top}\bm{L}\bm{f}}
       {\bm{f}^{\top}\bm{D}^{(w)}\bm{f}}
  =\frac46=\frac23.
\]
广义特征问题的第二特征值和一个对应方向则为
\[
  \lambda_{2,\mathrm{sym}}=\frac12,\qquad
  \bm{y}_2\propto(1,\,1/2,\,-1/2,\,-1)^{\top},
  \qquad
  \bm{L}\bm{y}_2=\frac12\bm{D}^{(w)}\bm{y}_2.
\]
阈值切分仍返回中间一刀，枚举也给出离散NCut最优值\(2/3\)。两个松弛在这个对称小图上
得到同一分割，但特征值和特征向量不同，正好说明RatioCut与Normalized Cut不能混成
同一个特征问题。

松弛向量不是离散分割本身。按符号或阈值切分会产生一个可行割，其目标值应回到原离散
目标计算。谱间隙小还意味着特征向量容易受扰动旋转；重复特征值时，稳定对象是特征子空间，
不是某个算法任意选出的单一基向量。

\subsection{扩散用矩阵函数抑制高频变化}

连续时间图扩散定义为
\begin{equation}
  \frac{\mathrm{d}\bm{f}(t)}{\mathrm{d}t}
  =-\bm{L}\bm{f}(t),
  \qquad
  \bm{f}(0)=\bm{f}_0.
  \label{eq:graph-theory-heat-equation}
\end{equation}
第~\ref{chap:matrix-analysis}章的矩阵指数给出
\[
  \bm{f}(t)=\exp(-t\bm{L})\bm{f}_0
  =\sum_{k=1}^{n}e^{-t\lambda_k}
  \langle\bm{u}_k,\bm{f}_0\rangle\bm{u}_k.
\]
特征值大的模式衰减更快，零特征值对应的分量保持不变。正权支撑图连通时，
\(\bm{f}(t)\)趋向初始信号的常量投影；支撑图不连通时，每个支撑连通分量分别趋向
自己的平均。

总和保持来自\(\bm{1}^{\top}\bm{L}=\bm{0}^{\top}\)：
\[
  \frac{\mathrm{d}}{\mathrm{d}t}
  \bm{1}^{\top}\bm{f}(t)=0.
\]
平方范数与Dirichlet能量是两个不同的量。平方范数的导数恰好是负的Dirichlet能量：
\[
  \frac{\mathrm{d}}{\mathrm{d}t}
  \frac12\lVert\bm{f}(t)\rVert_2^2
  =
  -\bm{f}(t)^{\top}\bm{L}\bm{f}(t)
  \leq0.
\]
Dirichlet能量本身也下降，但对应的导数多出一个拉普拉斯：
\[
  \frac{\mathrm{d}}{\mathrm{d}t}
  \left[\bm{f}(t)^{\top}\bm{L}\bm{f}(t)\right]
  =
  -2\bm{f}(t)^{\top}\bm{L}^2\bm{f}(t)
  =
  -2\lVert\bm{L}\bm{f}(t)\rVert_2^2
  \leq0.
\]
第一个等式说明非平滑分量消耗平方范数，第二个等式直接说明边上平方差总量不会增加。
扩散因此削弱顶点间差异，但不表示它保留任务所需的类别边界。扩散时间过长会把同一
正权支撑连通分量内所有差异抹平。

\subsection{离散Dirichlet能量与连续几何的对应有条件}

式~\eqref{eq:graph-theory-dirichlet-energy}是第
~\ref{chap:differential-geometry-and-symmetry}章连续Dirichlet能量
\[
  \int_{\mathcal{M}}
  \lVert\operatorname{grad}f\rVert_g^2\,\mathrm{d}V_g
\]
的离散对应：连续梯度衡量无穷小方向的变化，图边差分只比较被连接的样本对。两者形式
相似，却不能只靠形式宣布收敛。

若图来自点云，邻域半径决定哪些局部方向被采样，核权重决定不同距离的贡献，归一化决定
是否保留采样密度，边界又会改变极限条件。邻域过大可能跨越弯曲流形的不同折，过小则使
图断连。样本密度不均时，未校正的图拉普拉斯可能逼近带密度漂移的算子，而不是纯粹由
Riemann体积定义的Laplace--Beltrami算子。

因此，图谱只相对于已经选定的图和权重给出精确代数结论。它是否近似某个连续空间，
需要额外的采样、尺度与归一化条件；它是否保留预测语义，又需要任务证据。

\subsection{顶点重编号要求表示按置换同步变化}

顶点名称只是坐标标签。同一图经置换矩阵\(\bm{P}\)重编号后，
\[
  \bm{W}'=\bm{P}\bm{W}\bm{P}^{\top},
  \qquad
  \bm{D}^{(w)\prime}=\bm{P}\bm{D}^{(w)}\bm{P}^{\top},
  \qquad
  \bm{L}'=\bm{P}\bm{L}\bm{P}^{\top}.
  \label{eq:graph-theory-laplacian-relabel}
\]
若顶点信号也重编号为\(\bm{f}'=\bm{P}\bm{f}\)，则
\[
  \bm{L}'\bm{f}'=\bm{P}\bm{L}\bm{f}.
\]
图拉普拉斯作用因此对顶点置换等变：输入顶点怎样重排，输出按同样方式重排。

更一般地，图滤波器若为多项式
\begin{equation}
  h(\bm{L})=\sum_{k=0}^{K}a_k\bm{L}^k,
  \label{eq:graph-theory-polynomial-filter}
\end{equation}
则
\[
  h(\bm{L}')\bm{f}'
  =\sum_{k=0}^{K}a_k(\bm{P}\bm{L}\bm{P}^{\top})^k\bm{P}\bm{f}
  =\bm{P}h(\bm{L})\bm{f},
\]
因为\(\bm{P}^{\top}\bm{P}=\bm{I}\)。若再对顶点输出作求和等对称汇总，可以得到置换
不变的图级表示。等变性只保证结果不依赖任意编号，不保证不同非同构图一定得到不同表示。

若存在置换矩阵\(\bm{P}\)使
\(\bm{A}'=\bm{P}\bm{A}\bm{P}^{\top}\)，两张图称为
\term{同构}（isomorphic）。同构图在忽略顶点名称后具有相同连接结构。比较两个给定
矩阵是否仅差重编号，不能使用一般特征分解替代：同谱图可能不同构，重复特征值下特征向量
基也不唯一。

\subsection{多重集合细化刻画局部组合信息}

邻居没有天然顺序，所以一个顶点的局部结构应通过邻居标记的
\term{多重集合}（multiset）表示。多重集合既记录元素种类，也记录每种元素出现次数；
\(\{a,a,b\}\)与\(\{a,b\}\)不同，但不区分排列次序。

一维Weisfeiler--Leman细化（1-dimensional Weisfeiler--Leman refinement,
1-WL）从初始标记\(c_i^{(0)}\)出发，反复令
\begin{equation}
  c_i^{(t+1)}=\operatorname{HASH}\left(
    c_i^{(t)},
    \left\{\!\left\{
      c_j^{(t)}:j\in N(i)
    \right\}\!\right\}
  \right),
  \label{eq:graph-theory-wl-refinement}
\end{equation}
其中\(\operatorname{HASH}\)对不同输入对给出不同新标记。每轮都把顶点自身旧标记与
邻居标记多重集组合，因此同构重编号只会同步重排标记，不改变各标记的出现次数。

这个过程会在有限步后稳定。令\(\mathcal{P}^{(t)}\)表示按
\(c_i^{(t)}\)相等关系得到的顶点分区。由于更新输入包含顶点自身旧标记，旧标记不同的
两个顶点下一轮仍不可能获得相同标记，所以
\(\mathcal{P}^{(t+1)}\)只能细分\(\mathcal{P}^{(t)}\)，不能合并其中的块。若分区
尚未稳定，块数至少增加\(1\)；而\(n\)个顶点至多分成\(n\)个非空块。因此从任意初始
分区出发，严格细分至多发生\(n-1\)次。某轮分区不再改变后，下一轮每个块中的顶点仍有
相同的自身标记和邻居标记计数，后续分区也保持不变。这里的“稳定”指等价类不再细分，
哈希实现仍可为这些类换用新的符号名称。

细化稳定不表示已经解决图同构。若所有顶点初始同标记，六顶点环\(C_6\)与两个互不相连
的三角形都让每个顶点始终看到两个同标记邻居；1-WL无法区分二者，尽管一张图连通、另一张
不连通。加入唯一顶点标识会人为破坏所需的重编号对称性，也不能作为无条件补救。

1-WL提供的是局部组合表达的参照。使用共享顶点更新和对称邻域聚合的具体图网络能否达到、
低于或超过该区分能力，要在模型正文中根据聚合函数、初始特征和读出方式证明。本章只建立
多重集合、置换作用和细化过程，不把结构等变性升级成学习能力或泛化保证。

\section{本章小结}
\label{sec:graph-theory-summary}

图把有限变量间的直接关系从候选集合中单独提取出来。顶点和边规定关系骨架，有向、无向、
加权、二部与超图分别保留不同类型的信息；邻接矩阵把邻域汇总写成线性映射，但矩阵位置
依赖顶点编号。路径、环、连通分量、树与DAG描述局部边怎样组成整体结构，仍不自动附带
概率、因果或代价语义。

结构能够改变计算组织。分隔集把图切成通过小接口相连的部分，消元则揭示中间函数会在哪些
邻居间产生新联系。贯穿例中，先处理顶点\(1,4\)只需宽度\(2\)，先处理顶点\(2\)
却会补出边\(\{1,4\}\)并产生四顶点团；
树分解与树宽把这种差异变成局部袋大小。树上子树的Helly性质进一步保证每个团都完整
落入某个袋，从而使三角形等局部团给出严格的树宽下界。小树宽使精确动态规划的指数依赖
局部宽度，而不是全部变量数，但寻找好分解和处理大取值域仍有成本。

结构之外还需要代数。分配律使局部汇总可以移出无关因子，半环则统一总权重、最大权重、
最小成本与可达性计算。同一消元结构换用和积、最大积或最小加法会回答不同问题；求和与
最大化混用时不能任意交换次序，因为这会改变哪些选择能够依赖尚未汇总的变量。

图算法的精确性各有条件。带起点与不可达边界的Bellman式在DAG上可按拓扑序递推；
一般有环图需要反复松弛，并单独处理目标相关负环。非负边长支持Dijkstra的贪心确定，
割性质支持最小生成树，残量网络把最大流与最小割连接起来，二部结构把匹配转成整数流。
树宽小支持局部动态规划；二值成对能量只有满足次模不等式时，才可分解成非负割边和一元
终端边，由一次最小割精确求解。一般离散目标仍需松弛、取整或搜索，并用原目标值与上下界
评价结果。

图拉普拉斯从另一侧分析同一关系结构。它的二次型是边上平方差之和，零空间记录正权支撑图
的连通分量，低特征值方向描述沿强边缓慢变化的信号。RatioCut的松弛使用组合拉普拉斯的
费德勒向量，Normalized Cut则使用度加权广义特征问题，二者不能混用。矩阵指数按谱衰减
高频差异；这个离散能量能否近似第~\ref{chap:differential-geometry-and-symmetry}章
的连续Dirichlet能量，取决于图构造、权重、采样密度、尺度与边界。顶点重编号下，
拉普拉斯和多项式滤波按置换等变；1-WL进一步用邻居标记多重集刻画局部组合信息，但仍
可能无法区分某些非同构图。

由此，章首问题可以分成两层回答：连接和分解决定哪些局部计算能够复用以及中间对象多大，
半环、容量、代价或顶点信号则决定计算结果的含义。第~\ref{chap:signal-analysis}章将把
谱坐标用于有顺序的信号，比较规则序列频率与图频率，并把线性递推展开成随时间组织的状态核。
